\documentclass[preprint,11pt]{elsarticle}
\usepackage{amsmath,amssymb,amsthm,bm,mathtools}
\usepackage{booktabs,multirow,array}
\usepackage[table]{xcolor}
\usepackage{graphicx}
\usepackage{hyperref}
\newcommand{\todo}[1]{\textcolor{red}{[\textbf{TODO}: #1]}}
\newcommand{\todocite}{\textcolor{red}{[?]}}
\newcommand{\R}{\mathbb{R}}
\newcommand{\x}{\bm{x}}
\newcommand{\z}{\bm{z}}
\newcommand{\W}{\bm{W}}
\newcommand{\A}{\bm{A}}
\newcommand{\B}{\bm{B}}
\newcommand{\g}{\bm{g}}
\DeclareMathOperator{\diag}{diag}
\DeclareMathOperator{\rank}{rank}
\newtheorem{proposition}{Proposition}
\newtheorem{lemma}{Lemma}
\newtheorem{definition}{Definition}

\journal{Neurocomputing}

\begin{document}
\begin{frontmatter}
\title{Expressive Low-Rank Adaptation and the Evaluation Artifacts that Obscure It:\\ A Unified Study of Nonlinear and Input-Dependent LoRA}

\author{Hao Chen \todocite}
\address{\todo{affiliation}}

\begin{abstract}
Low-Rank Adaptation (LoRA) fine-tunes a frozen model with a \emph{linear, input-independent} rank-$r$ update. We study
whether two orthogonal forms of added expressiveness help: (i) a \emph{nonlinearity} inside the low-rank code (the LeNA
framework, instantiated with the published AuroRA activation), and (ii) an \emph{input-dependent} update (IQ-LoRA). We give
a single formulation that contains all three as special cases, exposing exactly one extra degree of freedom each. Empirically, on GSM8K reasoning with Llama-2-7B, both
variants \emph{match} well-tuned LoRA at an identical parameter budget across ranks $r\in\{1,2,4,16,32\}$. We then show
that the opposite conclusion---an apparent large advantage for the expressive variants---is produced by three concrete and
common evaluation artifacts: (a) \emph{answer-extraction bias} under trailing-token degeneration, (b) \emph{repetition
degeneration} of the strong baseline, and (c) \emph{template collapse} of perturbation benchmarks. We formalize each,
give an artifact-free protocol, and reproduce how a spurious ``$16\times$ parameter-efficiency'' claim collapses to parity
under it. The result is a disciplined map of when extra expressiveness does and does not help, and a reusable protocol for
comparing parameter-efficient fine-tuning (PEFT) methods on generative reasoning benchmarks.
\end{abstract}
\begin{keyword}
Parameter-efficient fine-tuning \sep Low-rank adaptation \sep LLM reasoning \sep Evaluation protocol
\end{keyword}
\end{frontmatter}

%==========================================================
\section{Introduction}
Adapting a large pretrained model to a task by updating a small number of parameters is now standard. LoRA
\todocite{} is the dominant choice: it freezes a weight $\W_0\in\R^{d_o\times d_i}$ and learns a rank-$r$ correction
$\Delta\W=\tfrac{\alpha}{r}\B\A$ with $\A\in\R^{r\times d_i}$, $\B\in\R^{d_o\times r}$, $r\ll\min(d_o,d_i)$. Two properties
define LoRA and bound its hypothesis class: the correction is \emph{linear} in the low-rank code $\z=\A\x$, and it is
\emph{static}---the same $\Delta\W$ for every input.

A natural question, and the subject of a growing literature, is whether relaxing either property helps. We isolate the two
relaxations:
\begin{itemize}
\item \textbf{Nonlinearity (LeNA).} Insert a learnable elementwise nonlinearity $\phi$ between the down- and up-projection,
so the code$\to$update map is no longer linear.
\item \textbf{Input-dependence (IQ-LoRA).} Let the update depend on the input, $\Delta\W(\x)$, through a second projection,
so the effective correction is instance-conditional.
\end{itemize}
Both are intuitive and both enlarge the function class strictly (Section~\ref{sec:method}). One expects a gain.

\paragraph{Findings.} Our controlled study says otherwise, and says it carefully.
(1)~At matched parameter budget the nonlinear axis (AuroRA) \emph{matches} well-tuned LoRA on GSM8K across every rank we
test (Table~\ref{tab:rankcurve}); the input-dependent axis (IQ-LoRA) is expected to match on the same evidence
(\textcolor{red}{Table~\ref{tab:datasets}, to run}). Neither degree of freedom yields a robust improvement.
(2)~The literature-style \emph{opposite} conclusion is recoverable, but only through evaluation artifacts, which we name
and formalize (Section~\ref{sec:protocol}): extraction bias, repetition degeneration, and template collapse. Under a
naive-but-common protocol the nonlinear variant appears to reach LoRA's rank-$32$ accuracy at rank $2$
(a ``$16\times$ efficiency'' claim); under an artifact-free protocol the same checkpoints are at parity
(Table~\ref{tab:artifact}). Measuring that parity correctly---not any new adapter---is the contribution.

\paragraph{Contributions.}
\begin{enumerate}
\item A \emph{unified formulation} (Section~\ref{sec:method}) in which LoRA, LeNA, and IQ-LoRA differ by exactly one term,
making the added degree of freedom explicit and the comparison strictly controlled.
\item A \emph{rigorous evaluation protocol} (Section~\ref{sec:protocol}) for generative reasoning benchmarks that is
provably invariant to the three named artifacts, together with a $\sim$9$\times$ faster batched implementation.
\item An \emph{honest empirical map} (Section~\ref{sec:exp}): expressive low-rank variants match LoRA at matched budget,
and prior apparent gains are attributable to the artifacts. Parity, correctly measured, is the result.
\end{enumerate}

%==========================================================
\section{Related Work}
\paragraph{Static low-rank adaptation.} LoRA \todocite{} learns a linear rank-$r$ update; its refinements keep the update
\emph{linear and input-independent} while improving where or how the rank budget is spent---DoRA \todocite{} decouples
magnitude and direction, AdaLoRA \todocite{} allocates rank by importance, and unified accounts \todocite{} place these,
adapters \todocite{} and IA$^3$ \todocite{} in one design space. Our two axes are orthogonal to this line: they change the
\emph{form} of the update (nonlinear; input-dependent) rather than its rank allocation.

\paragraph{Expressive adapters.} Nonlinearity inside a bottleneck adapter is classical \todocite{}, and recent work learns
the activation itself---rational/Kronecker maps and the AuroRA activation \citep{aurora} and LoRAN \citep{loran} that we
adopt as the \emph{nonlinear instances} studied here. Input-conditional adaptation appears as hypernetwork- and
conditional adapters \todocite{}; IQ-LoRA is a lightweight instance that makes the update quadratic in the input via a
single extra projection (Eq.~\ref{eq:iq-operator}). We do not propose these activations or conditioning schemes; we ask,
under a controlled budget and a fair protocol, whether either \emph{form} of expressiveness helps.

\paragraph{Evaluating LLM reasoning.} GSM8K \todocite{} and its perturbations---GSM-Symbolic \todocite{}, SVAMP/ASDiv
\todocite{}---score a decoded string. Prior studies note sensitivity to answer normalization and decoding \todocite{};
we make this quantitative for PEFT comparison, isolating three artifacts (extraction, degeneration, template collapse) and
giving an invariant protocol (Section~\ref{sec:protocol}).

%==========================================================
\section{Method: A Unified View of Expressive Low-Rank Adaptation}
\label{sec:method}
\paragraph{Base module.} A PEFT module acts on a frozen linear map. Writing the low-rank code $\z=\A\x\in\R^{r}$, LoRA is
\begin{equation}
h_{\text{LoRA}}(\x)=\W_0\x+\tfrac{\alpha}{r}\,\B\,\z,\qquad \z=\A\x.
\label{eq:lora}
\end{equation}
The map $\z\mapsto \tfrac{\alpha}{r}\B\z$ is linear, and $\A,\B$ are fixed at inference. We now add one degree of freedom at
a time; both reduce to \eqref{eq:lora} when their new parameter is switched off, giving a strictly controlled comparison.

\subsection{Nonlinearity axis: a nonlinear code}
The first axis inserts a learnable elementwise nonlinearity $\phi$ into the code, replacing the linear map $\z\mapsto\B\z$:
\begin{equation}
h_{\phi}(\x)=\W_0\x+\B\,\phi\!\big(\g\odot \nu(\A\x)\big),
\label{eq:lena}
\end{equation}
with normalization $\nu$ (RMS over the $r$ channels) and a learned per-channel gate $\g\in\R^{r}$. We instantiate $\phi$
with \textbf{AuroRA}, a \emph{published} nonlinear activation \citep{aurora}\footnote{arXiv:2505.18738.}, not a
contribution of this work:
\begin{equation}
\phi_{\text{AuroRA}}(\z)=\tanh\!\big(\bm{H}\tanh(\z)\big)+\bm{w}_s\odot \mathrm{spline}(\z),
\label{eq:aurora}
\end{equation}
where $\bm{H}\in\R^{\tilde r\times\tilde r}$ is a cross-channel map, $\bm{w}_s\in\R^{\tilde r}$ per-channel spline weights,
and $\mathrm{spline}(\cdot)$ a learned-knot piecewise interpolant. \eqref{eq:lena} recovers LoRA when $\phi=\mathrm{id}$,
$\g=\bm1$, $\nu=\mathrm{id}$; for nonlinear $\phi$ the map $\z\mapsto\B\phi(\cdot)$ is not any rank-$r$ linear map, so the
hypothesis class strictly contains LoRA's. \textcolor{red}{[Positioning: AuroRA and LoRAN (\citep{loran}, EMNLP'24) are
prior activations; this axis \emph{studies} them under our protocol, it does not claim them. If a genuinely novel
activation is to be centered, replace \eqref{eq:aurora} accordingly---pending author input.]}

\subsection{IQ-LoRA: input-dependence of the update}
IQ-LoRA multiplicatively modulates the code by a second projection of the input, making the update instance-conditional:
\begin{equation}
h_{\text{IQ}}(\x)=\W_0\x+\B\Big(\z\odot\big(\bm{1}+\lambda\,\tanh(\A_2\x)\big)\Big),\qquad \z=\A\x,
\label{eq:iq}
\end{equation}
with $\A_2\in\R^{r\times d_i}$ and scale $\lambda\ge0$. The bounded gate $\tanh$ keeps the modulation in $(1\!-\!\lambda,
1\!+\!\lambda)$ (variant IQG); the raw form uses $\z+\lambda(\z\odot\A_2\x)$. Equivalently the effective correction is an
\emph{input-dependent} low-rank operator
\begin{equation}
\Delta\W(\x)=\B\,\diag\!\big(\bm{1}+\lambda\tanh(\A_2\x)\big)\,\A,
\label{eq:iq-operator}
\end{equation}
which is quadratic in $\x$ and reduces to LoRA at $\lambda=0$.
\begin{proposition}[Input-dependence]
\label{prop:iq}
For $\lambda>0$ the input--output map $\x\mapsto\Delta\W(\x)\,\x$ of \eqref{eq:iq-operator} is not linear in $\x$, hence
is not realized by any fixed operator $\Delta\W$; the IQ-LoRA family is therefore not contained in the static rank-$r$
family. At $\lambda=0$ it equals \eqref{eq:lora}.
\end{proposition}
\begin{proof}
$\Delta\W(\x)\,\x=\B\z+\lambda\,\B\big(\z\odot\tanh(\A_2\x)\big)$ with $\z=\A\x$. The second summand contains the product
$(\A\x)\odot\tanh(\A_2\x)$, a non-affine function of $\x$; thus the map is nonlinear and no constant $\Delta\W$ reproduces
it. Setting $\lambda=0$ deletes the summand. \qed
\end{proof}
The extra cost is one projection $\A_2\in\R^{r\times d_i}$; in experiments we charge it against a LoRA control of matched
trainable count (Table~\ref{tab:params}).

\paragraph{Parameter accounting.} All three share the $r(d_i+d_o)$ cost of $(\A,\B)$. The nonlinearity axis adds only the
gate and activation parameters, $O(r+\tilde r^2)$, negligible against $r(d_i+d_o)$ for $\tilde r\le r\ll d$. The
input-dependence axis adds a full projection $\A_2$, $r\,d_i$; we therefore compare IQ-LoRA not to LoRA at the same $r$ but
to a LoRA control whose rank is enlarged to absorb the same trainable count (Table~\ref{tab:params}), so any difference is
attributable to the mechanism, not the budget.
\begin{table}[t]\centering
\caption{Trainable parameters per adapted projection ($d_i{=}d_o{=}d$). $\tilde r\!\le\! r$ is the AuroRA inner width;
$K$ spline knots. IQ-LoRA is budget-matched by an enlarged-rank LoRA control.}
\label{tab:params}
\begin{tabular}{lll}
\toprule
Method & Trainable parameters & Overhead vs.\ LoRA\\
\midrule
LoRA          & $2rd$ & --\\
LeNA/AuroRA   & $2rd+\underbrace{r}_{\g}+\underbrace{\tilde r^2+\tilde rK}_{\phi}$ & $O(r+\tilde r^2)$, negligible\\
IQ-LoRA       & $2rd+\underbrace{rd}_{\A_2}=3rd$ & $+rd$ (matched by LoRA at $r'{=}\tfrac{3}{2}r$)\\
\bottomrule
\end{tabular}
\end{table}

%==========================================================
\section{A Rigorous Evaluation Protocol}
\label{sec:protocol}
Generative reasoning benchmarks score a decoded string, not a label. The scoring pipeline---answer extraction, decoding,
and subset selection---has degrees of freedom that, we show, can dominate the method effect. We formalize three artifacts
and give an invariant protocol.

\subsection{Artifact 1: extraction bias under trailing degeneration}
GSM8K answers terminate with \texttt{\#\#\#\# $n$}. Let a decoded completion be $y$ and let $E:\;y\mapsto n$ be an
extraction operator. Two common choices:
\begin{align}
E_{\text{last}}(y)&=\text{last numeric token of }y, &
E_{\text{first}}(y)&=\text{number in the \emph{first} ``\texttt{\#\#\#\#}'' match of }y.
\end{align}
\begin{definition}[Degeneration-invariance]
$E$ is degeneration-invariant if $E(y)=E(y\Vert s)$ for every trailing string $s$ appended after the first terminal answer
in $y$. $E_{\text{first}}$ is degeneration-invariant; $E_{\text{last}}$ is not.
\end{definition}
\begin{lemma}[Non-invariant extraction has method-dependent bias]
\label{lem:bias}
Let $y$ carry a correct terminal answer and let $p_M$ be the probability that method $M$ appends a further numeric token
after it. Under $E_{\text{last}}$, every such continuation scores the example wrong, so the measured accuracy is
$\mathrm{acc}_M-p_M\!\cdot\!\Pr[\text{correct}]$ instead of $\mathrm{acc}_M$; the induced downward bias $p_M\Pr[\text{correct}]$
is method-dependent and can be arbitrarily large. $E_{\text{first}}$ removes it ($p_M$ drops out).
\end{lemma}
\begin{proof}
Appending any token $t\neq$ the answer makes $E_{\text{last}}(y\Vert t)\neq\text{gold}$ while $E_{\text{first}}(y\Vert t)=
\text{gold}$ by Definition~1; taking expectations over the $p_M$-fraction of correct-but-continued completions gives the
stated bias, which scales with $p_M$. \qed
\end{proof}
This is not academic: strong LoRA checkpoints frequently degenerate into ``\texttt{\#\#\#\# $n$ \#\#\#\# $n$}$\cdots$''
($15$--$36$ repetitions observed), giving a large $p_{\text{LoRA}}$ that $E_{\text{last}}$ converts into a spurious
baseline deficit---and hence a spurious advantage for any method that happens to degenerate less.

\subsection{Artifact 2: baseline-specific degeneration}
Lemma~\ref{lem:bias} bites only when $p_M$ differs across methods---and it does. LoRA checkpoints emit repeated
``\texttt{\#\#\#\#}'' blocks far more often than the expressive variants, so $E_{\text{last}}$ deflates the baseline
specifically, manufacturing the gap of Table~\ref{tab:artifact}. \todo{Figure: histogram of \#\#\#\# blocks per completion
(nHASH), LoRA vs AuroRA on GSM8K test; report $p_{\text{LoRA}}$ vs $p_{\text{AuroRA}}$ to instantiate Lemma~\ref{lem:bias}.}

\subsection{Artifact 3: template collapse in perturbation benchmarks}
GSM-Symbolic is $T{=}100$ templates $\times$ $I{=}50$ instances. Evaluating the first $N\!<\!I$ examples samples a
\emph{single} template---an estimate of one perturbation family reported as if it were the benchmark---whose mean and
variance need not track the full set. The invariant subset takes one instance per template. \todo{Table: mean$\pm$sd of
first-$N$ vs diverse-$N$ ($N{=}40,80$) for LoRA and AuroRA, quantifying the shift.}

\subsection{Batched implementation and a caveat}
$E_{\text{first}}$ with left-padded batched decoding is $\sim$9$\times$ faster than per-example decoding and leaves the
scores unchanged for the static adapters studied here. One caveat bounds the batching: adapters that carry a
\emph{cumulative causal state} over the sequence (e.g.\ test-time-training branches) are corrupted by left padding, whose
pad positions enter the running state; such adapters must be decoded unbatched. It does not affect LoRA, LeNA, or IQ-LoRA,
whose per-token maps are state-free.

%==========================================================
\section{Experiments}
\label{sec:exp}
\paragraph{Setup.} Llama-2-7B, GSM8K, targets \{q,k,v,up,down\}, $3$ epochs, batch $2$, lr $3\times10^{-4}$, $7$k train
examples. Accuracy = exact-match under $E_{\text{first}}$ on the GSM8K test set, batched. \todo{state exact $n$, seeds.}

\subsection{Nonlinearity matches LoRA across ranks}
\begin{table}[t]\centering
\caption{GSM8K exact-match vs.\ rank, artifact-free protocol ($E_{\text{first}}$, batched). LeNA and LoRA are at
parity across all ranks (mean of $n$ seeds). Larger differences than $\pm0.02$ do not appear.}
\label{tab:rankcurve}
\begin{tabular}{lccccc}
\toprule
rank $r$ & 1 & 2 & 4 & 16 & 32\\
\midrule
LoRA               & .269 & .276 & .283 & .295 & .294\\
AuroRA \citep{aurora} & .262 & .285 & .293 & .314 & .282\\
$\Delta$           & $-.007$ & $+.008$ & $+.010$ & $+.019$ & $-.012$\\
\bottomrule
\end{tabular}
\end{table}

\begin{table}[t]\centering
\caption{The ``efficiency win'' is an extraction artifact. At $r{=}2$, the naive $E_{\text{last}}$ undercounts LoRA and
manufactures a LeNA advantage; $E_{\text{first}}$ removes it. Same checkpoints.}
\label{tab:artifact}
\begin{tabular}{lccc}
\toprule
$r{=}2$ (GSM8K) & LoRA & AuroRA & apparent gap\\
\midrule
$E_{\text{last}}$ (naive)  & .194 & .257 & $+.063$ (``$16\times$'')\\
$E_{\text{first}}$ (ours)  & .276 & .285 & $+.008$ (parity)\\
\bottomrule
\end{tabular}
\end{table}

\begin{figure}[t]\centering
\fbox{\parbox{0.82\linewidth}{\centering\vspace{2mm}\textcolor{red}{\textbf{FIGURE~1 TO PRODUCE.} Accuracy vs.\ rank
$r\!\in\!\{1,2,4,8,16,32\}$, one line each for LoRA / AuroRA / IQ-LoRA, mean $\pm$ s.d.\ over $\ge3$ seeds, GSM8K,
$E_{\text{first}}$. Visualizes Table~\ref{tab:rankcurve}; all curves should overlap within error bars (the parity claim).}
\vspace{2mm}}}
\caption{Rank curve (visual of Table~\ref{tab:rankcurve}), to render from the seed runs.}
\label{fig:rankcurve}
\end{figure}

\subsection{Do the two axes help? (breadth)}
Table~\ref{tab:rankcurve} establishes parity for the nonlinear axis on GSM8K. Two gaps remain: the input-dependent axis has
no protocol-consistent numbers yet, and a single task cannot support a parity \emph{claim}. Tables~\ref{tab:iqrank}
and~\ref{tab:datasets} close them.

\begin{table}[t]\centering
\caption{\textcolor{red}{TO RUN (E0).} Input-dependent axis under the protocol: IQ-LoRA ($\lambda{=}0.5$) and the bounded
variant IQG vs.\ a budget-matched LoRA control ($r'{=}\tfrac32 r$, Table~\ref{tab:params}), GSM8K, $E_{\text{first}}$,
$\ge3$ seeds. Cells = exact-match; expected $\Delta\approx0$.}
\label{tab:iqrank}
\rowcolors{2}{red!6}{red!6}
\begin{tabular}{lccc}
\toprule
rank $r$ & LoRA ($r'$) & \textcolor{red}{IQ-LoRA} & \textcolor{red}{IQG}\\
\midrule
2  & .28 & \textcolor{red}{--} & \textcolor{red}{--}\\
8  & \textcolor{red}{--} & \textcolor{red}{--} & \textcolor{red}{--}\\
16 & .30 & \textcolor{red}{--} & \textcolor{red}{--}\\
\bottomrule
\end{tabular}
\end{table}

\begin{table}[t]\centering
\caption{\textcolor{red}{TO RUN (E1).} Generalization across datasets at matched budget, $E_{\text{first}}$, diverse
subsets, $\ge3$ seeds. Cells = exact-match (mean). Claim survives iff each expressive row tracks LoRA within noise on every
column.}
\label{tab:datasets}
\rowcolors{2}{red!6}{red!6}
\begin{tabular}{lccccc}
\toprule
Method & GSM8K & \textcolor{red}{GSM-Sym (div.)} & \textcolor{red}{SVAMP} & \textcolor{red}{ASDiv} & \textcolor{red}{GSM-Plus}\\
\midrule
LoRA        & .29 & \textcolor{red}{--} & \textcolor{red}{--} & \textcolor{red}{--} & \textcolor{red}{--}\\
AuroRA      & .29 & \textcolor{red}{--} & \textcolor{red}{--} & \textcolor{red}{--} & \textcolor{red}{--}\\
IQ-LoRA     & \textcolor{red}{--} & \textcolor{red}{--} & \textcolor{red}{--} & \textcolor{red}{--} & \textcolor{red}{--}\\
\bottomrule
\end{tabular}
\end{table}

\subsection{Comparison to the PEFT family}
\begin{table}[t]\centering
\caption{\textcolor{red}{TO RUN (E2).} Against the broader PEFT family at matched trainable budget, GSM8K + one shift
column, $E_{\text{first}}$, $\ge3$ seeds. Purpose: the expressive variants match not just vanilla LoRA but the family.}
\label{tab:baselines}
\rowcolors{2}{red!6}{red!6}
\begin{tabular}{lcc}
\toprule
Method & GSM8K & \textcolor{red}{GSM-Sym (div.)}\\
\midrule
LoRA                    & .29 & \textcolor{red}{--}\\
\textcolor{red}{DoRA}    & \textcolor{red}{--} & \textcolor{red}{--}\\
\textcolor{red}{AdaLoRA} & \textcolor{red}{--} & \textcolor{red}{--}\\
\textcolor{red}{Houlsby Adapter} & \textcolor{red}{--} & \textcolor{red}{--}\\
\textcolor{red}{IA$^3$}  & \textcolor{red}{--} & \textcolor{red}{--}\\
AuroRA / IQ-LoRA (ours-studied) & .29 & \textcolor{red}{--}\\
\bottomrule
\end{tabular}
\end{table}

\subsection{The apparent gains are evaluation artifacts}
Table~\ref{tab:artifact} already shows the $r{=}2$ flip. The next three items make the mechanism (Lemma~\ref{lem:bias}) and
its scope explicit; they are the paper's core evidence.

\begin{table}[t]\centering
\caption{\textcolor{red}{TO RUN (E3a).} Protocol-flip table---the punchline. Every (method,dataset) scored under BOTH
$E_{\text{last}}$ and $E_{\text{first}}$; last column marks whether the pairwise verdict vs.\ LoRA flips. Shows the method
effect is dominated by the protocol effect.}
\label{tab:flip}
\rowcolors{2}{red!6}{red!6}
\begin{tabular}{llccc}
\toprule
Method & Dataset & $E_{\text{last}}$ & $E_{\text{first}}$ & verdict flips?\\
\midrule
AuroRA & GSM8K & .257 & .285 & \textcolor{red}{--}\\
IQ-LoRA & GSM8K & \textcolor{red}{--} & \textcolor{red}{--} & \textcolor{red}{--}\\
AuroRA & \textcolor{red}{GSM-Sym} & \textcolor{red}{--} & \textcolor{red}{--} & \textcolor{red}{--}\\
\multicolumn{5}{c}{\textcolor{red}{$\cdots$ all method$\times$dataset cells $\cdots$}}\\
\bottomrule
\end{tabular}
\end{table}

\begin{figure}[t]\centering
\fbox{\parbox{0.82\linewidth}{\centering\vspace{2mm}\textcolor{red}{\textbf{FIGURE~2 TO PRODUCE (E3b).} nHASH histogram:
distribution of the number of ``\texttt{\#\#\#\#}'' blocks per completion, LoRA vs.\ AuroRA, GSM8K test. Report
$p_{\text{LoRA}},p_{\text{AuroRA}}$ (fraction with $>\!1$ block). Instantiates $p_M$ of Lemma~\ref{lem:bias}: LoRA's mass at
$\ge2$ is the mechanism of the artifact.}\vspace{2mm}}}
\caption{Baseline-specific degeneration (evidence for Lemma~\ref{lem:bias}).}
\label{fig:nhash}
\end{figure}

\begin{table}[t]\centering
\caption{\textcolor{red}{TO RUN (E3c).} Template collapse (Artifact~3): GSM-Symbolic accuracy under the first-$N$ subset
vs.\ the one-per-template subset, mean$\pm$s.d.\ over seeds, $N\!\in\!\{40,80\}$. A large gap/variance shift means first-$N$
does not estimate the benchmark.}
\label{tab:template}
\rowcolors{2}{red!6}{red!6}
\begin{tabular}{lcc}
\toprule
subset & LoRA (mean$\pm$sd) & \textcolor{red}{AuroRA (mean$\pm$sd)}\\
\midrule
\textcolor{red}{first-40}   & \textcolor{red}{--} & \textcolor{red}{--}\\
\textcolor{red}{diverse-40} & \textcolor{red}{--} & \textcolor{red}{--}\\
\textcolor{red}{diverse-80} & \textcolor{red}{--} & \textcolor{red}{--}\\
\bottomrule
\end{tabular}
\end{table}

\begin{figure}[t]\centering
\fbox{\parbox{0.82\linewidth}{\centering\vspace{2mm}\textcolor{red}{\textbf{FIGURE~3 TO PRODUCE.} The money figure: for
each method, the apparent gap vs.\ LoRA under $E_{\text{last}}$ (red bars) vs.\ $E_{\text{first}}$ (grey bars), across
datasets. Red bars are large and inconsistent; grey bars cluster at zero---artifacts manufacture gaps that fair evaluation
erases.}\vspace{2mm}}}
\caption{Apparent method gap under naive vs.\ artifact-free protocols.}
\label{fig:money}
\end{figure}

\subsection{Ablations}
\begin{table}[t]\centering
\caption{\textcolor{red}{TO RUN (E4).} What the nonlinearity buys: fix the LeNA framework, vary $\phi\in$
\{identity(=LoRA), GELU, spline, AuroRA\} at $r{=}2,8$, GSM8K, $E_{\text{first}}$. Isolates whether any activation (not just
AuroRA) helps, and at which rank the bottleneck binds.}
\label{tab:ablate-act}
\rowcolors{2}{red!6}{red!6}
\begin{tabular}{lcc}
\toprule
$\phi$ & $r{=}2$ & $r{=}8$\\
\midrule
identity ($=$LoRA) & .276 & \textcolor{red}{--}\\
\textcolor{red}{GELU} & \textcolor{red}{--} & \textcolor{red}{--}\\
\textcolor{red}{spline} & \textcolor{red}{--} & \textcolor{red}{--}\\
AuroRA & .285 & \textcolor{red}{--}\\
\bottomrule
\end{tabular}
\end{table}

\begin{table}[t]\centering
\caption{\textcolor{red}{TO RUN (E5).} IQ-LoRA sensitivity: modulation scale $\lambda\in\{0,0.25,0.5,1\}$ and raw vs.\
bounded (IQG) gate, $r{=}8$, GSM8K, $E_{\text{first}}$. $\lambda{=}0$ is the LoRA anchor.}
\label{tab:ablate-iq}
\rowcolors{2}{red!6}{red!6}
\begin{tabular}{lcccc}
\toprule
$\lambda$ & 0 & 0.25 & 0.5 & 1\\
\midrule
\textcolor{red}{IQ (raw)} & .29 & \textcolor{red}{--} & \textcolor{red}{--} & \textcolor{red}{--}\\
\textcolor{red}{IQG (tanh)} & .29 & \textcolor{red}{--} & \textcolor{red}{--} & \textcolor{red}{--}\\
\bottomrule
\end{tabular}
\end{table}

%==========================================================
\section{Conclusion}
Along two natural axes---nonlinearity and input-dependence---extra expressiveness does not beat well-tuned LoRA on GSM8K at
matched budget. The seemingly contrary evidence is an evaluation artifact, and we give a protocol, with a proof
(Lemma~\ref{lem:bias}), that removes it. We claim generality only for the \emph{protocol}; whether the parity verdict holds
beyond the datasets and model we test is an empirical question our RED tables are designed to answer. The practical
takeaway is narrow but load-bearing: PEFT comparisons on generative reasoning benchmarks should report the extraction
operator and subset construction, because---as Lemma~\ref{lem:bias} shows---the method effect can be smaller than the
protocol effect.

\bibliographystyle{elsarticle-num}
%\bibliography{refs}
\end{document}
